% Imporditud: tools/import_eksperimendid.py
% Allikas: Eesti füüsikaolümpiaadi eksperimendi ülesannete
%   kogu (Rael Kalda, 2023), ülesanne Ü94, lahendus L94.
\setAuthor{EFO žürii}
\setRound{piirkonnavoor}
\setYear{2006}
\setNumber{G E1}
\setDifficulty{3}
\setTopic{geomeetriline optika}
\setVahendid{kumerlääts, nõguslääts, mõõtejoonlaud, taskulambipirn kõrgel alusel, lapik taskulambi patarei, 2 juhet, ekraan, valge paberileht}
\setPunktid{10}
\setAllikas{Eksperimendi kogu Ü94, lahendus lk 73}
\setLahendusseis{toores}

\prob{Läätse fookuskaugus}
Määrake nõgusläätse fookuskaugus.

\hint

\solu
Esimene lahendus: Põlev lamp asetatakse kumerläätsest võimalikult kaugele, et läätsele langeva valgusvihu võiks lugeda paralleelseks (piisav vahemaa 50 cm). Mõõdetakse kumerläätse fookuskaugus. Kumerläätse ette asetatakse nõguslääts. Kui nõgusläätse ja kumerläätse fookused ühilduvad, on pärast kumerläätse valgusvihk paralleelne. Paralleelsuse määramiseks asetatakse ekraanile paberileht ja märgitakse sellele valguslaigu läbimõõt. Ekraani nihutamisel peab valguslaik ekraanil jääma sama suureks. Mõõdetakse läätsede optiliste keskpunktide kaugused l. Nõgusläätse fookuskaugus fn = fk $-$ l. Idee (3 p). Kumerläätse fookuskauguse määramine (1 p). Nõgusläätse õigesse kohta paigaldamine (1 p). Valgusvihu paralleelsuse kontroll (1 p). Läätsede vahekauguse mõõtmine (1 p). Nõgusläätse fookuskauguse arvutamine (1 p). Korduvkatsed (1 p). Veaarvutus (1 p).

Teine lahendus: Põlev lamp asetatakse kumerläätsest võimalikult kaugele, et läätsele langeva valgusvihu võiks lugeda paralleelseks. Mõõdetakse kumerläätse fookuskaugus. Kumerläätse taha asetatakse nõguslääts. Kui kumerläätse ja nõgusläätse fookused ühilduvad, on pärast nõgusläätse valgusvihk paralleelne. Kontrollitakse valgusvihu paralleelsus. Mõõdetakse läätsede optiliste keskpunktide kaugused l. Nõgusläätse fookuskaugus fn = fk $-$ l. Idee (3 p). Kumerläätse fookuskauguse määramine (1 p). Nõgusläätse õigesse kohta paigaldamine (1 p). Valgusvihu paralleelsuse kontroll (1 p). Läätsede vahekauguse mõõtmine (1 p). Nõgusläätse fookuskauguse arvutamine (1 p). Korduvkatsed (1 p). Veaarvutus (1 p).

Kolmas lahendus: Kui valgusallikas asub piisavalt kaugel, siis võib lugeda, et läätsele langevad kiired on paralleelsed. Nõguslääts asetatakse valgusallika ja ekraani vahele. Nõguslääts eemaldatakse ekraanist. Teatud kaugusel ekraanist tekib läätse võru ümber hele rõngas. Mõõdetakse läätse läbimõõt, heleda rõnga läbimõõt ja läätse kaugus ekraanist ning sarnastest kolmnurkadest arvutatakse nõgusläätse fookuskaugus. Sarnastest kolmnurkadest ilmneb, et D/d = a/f , millest f = ad/D. Meetod on eelnevatest ebatäpsem, sest ruumi valgustatus segab katse läbiviimist ning heleda rõnga piire ei ole võimalik kuigi täpselt fikseerida, seepärast on antud lahenduse eest maksimaalselt võimalik saada (6 p).

Neljas lahendus: Põlev lamp asetatakse kumerläätse fookusesse. Pärast läätse on valgusvihk paralleelne. Nõguslääts asetatakse heledasse paralleelsesse valgusvihku. Nõgusläätse taha paigutatakse ekraan. Nõguslääts eemaldatakse ekraanist. Teatud kaugusel ekraanist tekib läätse võru ümber hele rõngas. Mõõdetakse läätse läbimõõt, heleda rõnga läbimõõt ja läätse kaugus ekraanist ning sarnastest kolmnurkadest arvutatakse nõgusläätse fookuskaugus. Sarnastest kolmnurkadest ilmneb, et D/d = a/f , millest f = ad/D. Idee (3 p). Valgusallika paigutamine kumerläätse fookuskaugusesse (1 p). Valgusvihu paralleelsuse kontroll (1 p). Nõgusläätse läbimõõdu mõõtine (1 p). Ekraanil tekkinud valguslaigu läbimõõdu mõõtmine (1 p). Ekraani ja läätse vahelise kauguse mõõtmine (1 p). Korduvkatsed (1 p). Veaarvutus (1 p).
\probend
